\section{Related Work}
\label{sec:related-work}

\paragraph{Autonomous GUI exploration and environment models.}
GUI-explorer, UIExplore-Bench, and GUI-Bee treat interaction as a means to
discover environment-specific functionality or collect grounding data
\citep{xie2025guiexplorer,nica2025uiexplore,fan2025guibee}. UI-KOBE, EAM,
GraphPilot, and ActionEngine further organize experience into reusable graphs
or state-machine memories \citep{chai2026uikobe,qin2026eam,yu2026graphpilot,
zhong2026actionengine}. These systems establish exploration, structured memory,
and revisitation as important mechanisms. VERA changes the evidential status of
their typical outputs: a discovered function remains a hypothesis until an
observable outcome supports its admission to persistent knowledge.

\paragraph{Action-effect verification.}
VeriGUI predicts action effects and checks subsequent observations for
execution recovery, while VAGEN equips a verifier with interaction tools for
collecting task-completion evidence \citep{zhang2026verigui,cui2026vagen}.
GraphPilot and ActionEngine also validate or repair execution against stored
transitions \citep{yu2026graphpilot,zhong2026actionengine}. VERA uses the same
general insight---execution claims require environmental evidence---but assigns
verification a different decision role. In task-free model induction, the
verifier controls whether a location-conditioned functional claim becomes
reusable knowledge, rather than only whether the current task should continue
or recover.

\paragraph{Risk-aware GUI agents.}
WebGuard predicts outcomes and risk levels for state-changing Web actions
\citep{zheng2025webguard}; OS-Sentinel and OSGuard combine contextual judgment
with safety specifications or state invariants
\citep{sun2026ossentinel,mohammadmirzaei2026osguard}; and SeerGuard evaluates
proposed actions in the current GUI state before execution
\citep{yu2026seerguard}. VERA studies contextual risk at a different interface:
the assessed action is being used to acquire evidence for an unknown function.
Its risk record is therefore linked to the same claim, execution trace, and
outcome evidence that govern knowledge admission.

\paragraph{Positioning.}
VERA connects these directions around a single question: how should an
open-ended agent acquire reusable functional knowledge when both the conclusion
and the act of testing it carry risk? This yields a common trace from a frozen
hypothesis and pre-action risk record through execution evidence to an
admission decision. The contribution is this role assignment and connection in
task-free functional model induction, rather than a new foundation model,
browser executor, exploration primitive, or standalone risk taxonomy.
